\documentclass[11pt,a4paper]{article}

\usepackage[T1]{fontenc}
\usepackage[utf8]{inputenc}
\usepackage{mathpazo}
\usepackage[a4paper,margin=26mm,top=28mm,bottom=28mm]{geometry}
\usepackage{amsmath,amssymb}
\usepackage{booktabs}
\usepackage{graphicx}
\usepackage{microtype}
\usepackage{titlesec}
\usepackage{enumitem}
\usepackage{caption}
\usepackage{xcolor}
\usepackage{array}
\usepackage{tabularx}
\usepackage{float}
\usepackage{fancyhdr}
\usepackage[hidelinks]{hyperref}

\newcommand{\sectionbreak}{\clearpage}

% Figures are placed with [H] and sit where they are written. Every section
% starts a new page, so a deferred float would otherwise be flushed onto a
% page of its own and centred vertically with nothing around it.
\renewcommand{\topfraction}{0.9}
\renewcommand{\bottomfraction}{0.7}
\renewcommand{\textfraction}{0.1}
\renewcommand{\floatpagefraction}{0.8}
\setlength{\intextsep}{16pt plus 3pt minus 3pt}

\titleformat{\section}{\Large\bfseries}{\thesection}{0.8em}{}
\titlespacing*{\section}{0pt}{0pt}{18pt}
\titleformat{\subsection}{\large\bfseries}{\thesubsection}{0.7em}{}
\titlespacing*{\subsection}{0pt}{20pt}{8pt}
\titleformat{\subsubsection}{\normalsize\bfseries}{\thesubsubsection}{0.6em}{}
\titlespacing*{\subsubsection}{0pt}{14pt}{6pt}

\captionsetup{font=small,labelfont=bf,skip=8pt}
\setlength{\parskip}{6pt}
\setlength{\parindent}{0pt}
\linespread{1.06}

\pagestyle{fancy}
\fancyhf{}
\renewcommand{\headrulewidth}{0pt}
\fancyfoot[C]{\small\thepage}

\newcommand{\AU}{\mathrm{AU}}
\newcommand{\UJ}{\mathrm{UJ}}
\newcommand{\AJ}{\mathrm{AJ}}
\newcommand{\code}[1]{\texttt{\small #1}}

\begin{document}

% =====================================================================
\thispagestyle{empty}
\vspace*{2.2cm}

{\Huge\bfseries The AUD/JPY triangular basis\par}
\vspace{10pt}
{\Large Estimation and reconciliation at one-second resolution\par}

\vspace{16pt}
{\large AUD/JPY against AUD/USD $\times$ USD/JPY\\
1 July to 30 August 2024\par}

\vspace{2.2cm}
\hrule
\vspace{16pt}

\textbf{Scope.} The triangular basis is the residual of the parity relation between
AUD/JPY, AUD/USD and USD/JPY. This report documents the estimators used to measure that
residual at one-second resolution over a window containing the August 2024 yen carry
unwind, and reconciles the estimates that disagree. The residual remains inside the
transaction-cost band throughout the sample; no arbitrage is identified and none is
claimed.

\textbf{Principal result.} Basis dispersion over the stress window is $3.09$ times its
level in the calm segments when every second is pooled, and $1.22$ times when measured
within the hour. Section~\ref{sec:reconcile} shows that the difference is a variance
decomposition rather than a modelling disagreement. Day-to-day variation in the basis level
accounts for $80.6\%$ of the stress window's level variance, against $11.6\%$ in the calm
segments, and that split holds across every candidate regime boundary. Two days in that window fail the daily-bias check of
Section~\ref{sec:grid} by a factor of fifteen and account for the entire mean level shift
attributed to the regime.

\textbf{Results independent of that decomposition.} Hourly dispersion is best described by
a bounded excursion rather than a permanent shift, with posterior mass $1.000$ on the
two-changepoint models and an observed log Bayes factor of $+97.1$ against a null 95th
percentile of $-2.8$ (Section~\ref{sec:cp}). The transaction-cost band implied by quoted
spreads exceeds the median basis by a factor of $12.7$ at its tightest
(Section~\ref{sec:band}). The two dollar legs open the observed dislocations and the cross
leg closes them, established both by arithmetic decomposition and under a sign restriction
fixed before estimation (Sections~\ref{sec:vecm} and \ref{sec:attrib}).

\vspace{16pt}
\hrule
\vspace{20pt}

\textit{The emphasis throughout is on how each quantity is defined and why it is identified.
Every number cited is reproduced from a table written by the pipeline. The appendix maps
each claim to its file.}

\clearpage

% =====================================================================
\section{The estimand}

\subsection{The parity relation}

Let $S^{\AU}_t$, $S^{\UJ}_t$ and $S^{\AJ}_t$ denote mid quotes for AUD/USD, USD/JPY and
AUD/JPY. Absent frictions the cross rate is redundant, $S^{\AJ}_t = S^{\AU}_t S^{\UJ}_t$.
The object of study is the residual of that relation, expressed in AUD/JPY pips:
\begin{equation}
z_t \;=\; \frac{S^{\AJ}_t - S^{\AU}_t S^{\UJ}_t}{10^{-2}} .
\label{eq:basis}
\end{equation}

For the dynamic work the same quantity is written in logs. With
$x_t = (x^{\AU}_t, x^{\UJ}_t, x^{\AJ}_t)^{\top}$ the vector of log prices scaled by $10^4$
and
\begin{equation}
w = (-1,\,-1,\,+1)^{\top}, \qquad \tilde z_t = w^{\top} x_t ,
\label{eq:logbasis}
\end{equation}
$\tilde z_t$ agrees with \eqref{eq:basis} to first order. The ratio between the two scales
is measured at load rather than assumed, and holds to about two percent at the sample
average of the cross.

\subsection{The cointegrating vector is known}

Triangular parity is a cointegrating restriction on $x_t$ with a known vector. The rank is
one and the vector is $w$, both by arithmetic. Nothing about the relation is recovered from
the data, which removes the Johansen procedure and the question of how many cointegrating
relations the sample supports. Every dynamic statement below is therefore conditional on a
restriction that cannot be misspecified, and any failure to find mean reversion is a
statement about the market rather than about the specification.

\subsection{Three distinct estimands}

The phrase ``the basis widened'' names at least three quantities. They are defined here and
returned to in Section~\ref{sec:reconcile}.

\begin{enumerate}[leftmargin=2.2em,itemsep=6pt]
  \item \textbf{Within-hour dispersion.} $\operatorname{MAD}(z_t : t \in h)$ for an hour
  $h$. This is the feature the changepoint model of Section~\ref{sec:cp} fits.

  \item \textbf{Within-day dispersion.} $\mathbb{E}_d[\operatorname{Var}(z \mid d)]^{1/2}$,
  which pools hours within a day but not days.

  \item \textbf{Pooled dispersion.} $\operatorname{Var}(z)^{1/2}$ over every second in a
  regime, which additionally absorbs the movement of the daily mean.
\end{enumerate}

The three coincide only if the daily mean basis is constant. It is not, and under stress it
is far from constant. That fact accounts for most of the disagreement between sections of
this project.

\subsection{Episodes}

Dislocations in this sample are plateaus rather than isolated spikes, so several statements
are made at the level of an episode: a maximal run of consecutive seconds with $|z_t| > u$,
where $u = 2.42$ pips is the $99.95\%$ quantile of $|z|$ over July. Runs are split wherever
consecutive observations are more than one grid step apart, which prevents an episode from
spanning a market closure. Without that guard, exceedances either side of a weekend merge
into a single 48-hour episode. Fifty-eight episodes are detected.

% =====================================================================
\section{Constructing the observation grid}
\label{sec:grid}

\subsection{The failure a tick audit cannot detect}

The ingest stage audits every tick that exists: nulls, zeros, negatives, crossed quotes.
That audit is silent on ticks that are absent, because a missing tick leaves no row to
inspect. Forward-filling then substitutes a stale price, which is correct when a quote is
late and wrong when a feed is dead. Nothing in the filled series distinguishes the two
cases.

Two design decisions follow. Every row carries an explicit freshness flag per leg, taken
from whether that second contained a genuine tick rather than inferred from whether the
price changed, since a repeated identical quote is fresh and a difference cannot detect
that. And a separate gate runs between ingest and analysis.

\subsection{Two statistics from the over-identification}

Each leg is implied by the other two, for instance $\widehat{S^{\UJ}} = S^{\AJ}/S^{\AU}$,
and the implied series must agree with the quoted one. Two statistics are taken from that
gap. They detect different failures and neither detects the other's.

\begin{itemize}[leftmargin=1.6em,itemsep=5pt]
  \item The $99.9$th percentile of the \emph{absolute} gap, a tail statistic sensitive to a
  frozen or wildly wrong leg. Observed: $3.21$, $7.29$ and $4.83$ pips against a threshold
  of $25$.

  \item The \emph{median signed} gap within each day. A leg biased by a fraction of a pip
  moves this immediately and does not move the tail statistic at all, because the tail is
  set by ordinary noise an order of magnitude larger.
\end{itemize}

The second failure is the more consequential, because everything downstream continues to
look plausible and any analysis that differences, measures dispersion, or centres per day
removes the bias before it becomes visible. Section~\ref{sec:limits} reports what this check
returned.

Closures are excluded from both statistics. Across a weekend the grid is forward-filled
padding, so roughly $172{,}000$ identical values carry an identical frozen gap. Computed on
the full grid, the $99$th and $99.9$th percentiles were bit-identical for all three legs and
both equal to a single stale value.

\subsection{Closures identified from the data}

A second is classified as closed if it sits inside a run of at least ten minutes in which
none of the three pairs printed a new price. Ten minutes of simultaneous silence across
three majors does not occur while the market is open. Eight closures are found, all $48$
hours to within two minutes, and no trading calendar is consulted. About $3.86$ million
open-market seconds remain.

\subsection{Two measurement decisions}

\subsubsection*{The clock}

Vendor timestamps are documented as EST, but the sample lies entirely inside US daylight
saving, so the operative offset is UTC$-4$. The Friday-to-Sunday week boundary does not
establish this, since FX closes at 17:00 New York local year-round and that boundary appears
under either offset. The offset was determined by scanning candidate offsets against
UTC-stamped ticks from a second vendor over a control window. At $-4$h the correlation is
$1.00000$ with a gap standard deviation of $0.01$ pips; at $-5$h it is $-0.71$ and $48.65$
pips.

\subsubsection*{Quote pairing}

Bid and ask are carried through the grid beside the mid and are always taken from a single
tick. Three \code{ARG\_MAX} calls over one ordering key resolve to the same winning row, so
the bid and ask columns are null in identical rows and the forward fill cannot pair a bid
from one second with an ask from another. Taking $\max(\text{ask})$ with $\min(\text{bid})$
across a second would produce a tighter aggregate and a spread no participant was shown.
Section~\ref{sec:band} depends on this pairing, so it is asserted after the splice rather
than assumed.

% =====================================================================
\section{Exact Bayesian segmentation}
\label{sec:cp}

\subsection{The feature}

For each retained hour $h$, dispersion is summarised by the median absolute deviation
rescaled to standard-deviation units:
\[
\operatorname{MAD}_h \;=\; 1.4826 \cdot \operatorname{median}_{t \in h}
\bigl| z_t - \operatorname{median}_{s \in h} z_s \bigr| .
\]
The basis has excess kurtosis of $277$ over the full sample, so the hourly standard
deviation is determined by a small number of seconds while the MAD is not. The model is
fitted to $\log \operatorname{MAD}_h$, which makes a shift multiplicative and makes the
feature close to symmetric over the roughly $3600$ observations per hour.

Two nuisance profiles are subtracted. With $\mathrm{slot}(h)$ the within-day position of the
bucket and $\mathrm{age}(h)$ the buckets elapsed since the session reopened,
\begin{equation}
y_h \;=\; \log \operatorname{MAD}_h \;-\; m\bigl(\mathrm{slot}(h)\bigr)
\;-\; r\bigl(\min(\mathrm{age}(h), R)\bigr),
\label{eq:deseason}
\end{equation}
where $m$ and $r$ are medians taken over the full sample. The diurnal amplitude is roughly
fourfold peak to trough. The reopen term is required because a weekend closure is followed
by a thin Sunday evening, so the first buckets of each week sit systematically high; without
it, a regime ending during a closure is recorded a few hours after the reopen, since no
bucket exists inside the gap for a changepoint to occupy.

Both profiles are estimated over the whole sample and therefore cannot have been fitted
either side of a candidate date. A common profile cannot create a level shift, only reduce
one, so the adjustment works against the finding.

\subsection{Marginal likelihood of a segment}

Write a segment as a one-regressor regression
\[
\zeta_t \;=\; \mu \, \omega_t + e_t, \qquad e_t \sim \mathcal{N}(0, \sigma^2),
\]
which reduces to the ordinary segment mean when $\omega_t \equiv 1$ and also accommodates
the AR(1) transform below without modification. Place a Normal-Inverse-Gamma prior,
\[
\mu \mid \sigma^2 \sim \mathcal{N}\!\left(m_0, \tfrac{\sigma^2}{\kappa_0}\right),
\qquad
\sigma^2 \sim \mathrm{Inv\text{-}Gamma}(a_0, b_0),
\]
and let $n$, $S_{\omega\omega} = \sum \omega_t^2$, $S_{\omega\zeta} = \sum \omega_t \zeta_t$
and $S_{\zeta\zeta} = \sum \zeta_t^2$ be the sufficient statistics. Integrating out
$(\mu,\sigma^2)$ gives a closed form. With
\[
\kappa_n = \kappa_0 + S_{\omega\omega}, \quad
a_n = a_0 + \tfrac{n}{2}, \quad
m_n = \frac{\kappa_0 m_0 + S_{\omega\zeta}}{\kappa_n}, \quad
b_n = b_0 + \tfrac{1}{2}\!\left(S_{\zeta\zeta} + \kappa_0 m_0^2 - \kappa_n m_n^2\right),
\]
the log marginal likelihood is
\begin{equation}
\log p(\text{segment}) \;=\;
\log\Gamma(a_n) - \log\Gamma(a_0)
+ a_0 \log b_0 - a_n \log b_n
+ \tfrac{1}{2}\bigl(\log \kappa_0 - \log \kappa_n\bigr)
- \tfrac{n}{2}\log 2\pi .
\label{eq:nig}
\end{equation}

Every input to \eqref{eq:nig} is available in $O(1)$ from cumulative sums, which is what
makes exhaustive enumeration cheaper than sampling.

The prior is weak but proper. $\kappa_0 = 0.01$ is one hundredth of an hour's worth of
information about the segment mean, and $a_0 = 2$ is the smallest shape giving a finite
prior mean for $\sigma^2$. Propriety matters because improper priors make the model
comparison below undefined. The scale $b_0$ is set from the dispersion of the series being
fitted, which is empirical Bayes; that dispersion is inflated by the shift under test, so
the choice reduces the evidence for a change.

\subsection{AR(1) errors across a market that closes}

Hourly dispersion is strongly autocorrelated because volatility clusters. Under independent
errors a run of above-average hours is not distinguishable from a level shift, and the
estimator then selects an excursion on series containing no changepoint. Errors within a
segment are therefore modelled as AR(1) with coefficient $\rho$.

Whitening requires two cases, because the market closes weekly. Where bucket $h$ is adjacent
in clock time to $h-1$, the one-step-ahead form applies:
\[
\zeta_h = y_h - \rho\, y_{h-1}, \qquad \omega_h = 1 - \rho .
\]
Where it is not, meaning the first bucket of the sample and the first bucket after every
closure, there is no predecessor and the stationary marginal is used:
\[
\zeta_h = y_h \sqrt{1-\rho^2}, \qquad \omega_h = \sqrt{1-\rho^2}.
\]
Both place the observation on the same $\sigma$, so \eqref{eq:nig} applies unchanged. The
stationary form is also used at the first bucket of every segment, whose predecessor belongs
to a different segment with a different mean.

$\rho$ is marginalised over a grid $\mathcal{R}$ spanning $[0, 0.90]$ under a uniform prior
rather than profiled, so the reported evidence accounts for not knowing it:
\[
\log Z_M \;=\; \log \sum_{\rho \in \mathcal{R}} \exp\bigl(\log Z_M(\rho)\bigr)
- \log|\mathcal{R}| .
\]
The grid includes zero, so the independent-error model is nested and the data can return it.
The upper end stops short of one, where a level shift and a random walk are not separately
identified.

\subsection{Model space and enumeration}

Four models are compared over $n$ retained buckets, with $m$ the minimum admissible segment
length:
\begin{center}
\begin{tabular}{@{}llp{7.4cm}@{}}
\toprule
$M_0$ & no change & one segment $[0,n)$ \\
$M_1$ & permanent shift & one changepoint $\tau$ \\
$M_2$ & excursion & two changepoints; the outer segments share one set of parameters \\
$M_3$ & two free changes & two changepoints; three free segments \\
\bottomrule
\end{tabular}
\end{center}

$M_0$ is included because a one-changepoint model always returns a changepoint. Without a
no-change alternative in the comparison, the selected date carries no information about
whether a change occurred. $M_2$ separates the two substantive hypotheses: a new regime
against a bounded window followed by a return.

Changepoints are discrete and there are at most two, so the posterior is computed by
enumerating every admissible configuration. Writing
$\mathcal{T} = \{(\tau_1,\tau_2) : \tau_1 \geq m,\ \tau_2 - \tau_1 \geq m,\ n - \tau_2 \geq m\}$
and $\ell$ for the sum of \eqref{eq:nig} over the segments a configuration induces, the
evidence under a uniform prior over configurations is
\begin{equation}
\log Z_{M_k} \;=\; \log \sum_{c \in \mathcal{C}_k} \exp\bigl(\ell(c)\bigr)
\;-\; \log |\mathcal{C}_k| .
\label{eq:enumerate}
\end{equation}
The result is exact. There are no chains, no convergence diagnostics, and the output is
identical on every run. Posterior model probabilities follow from \eqref{eq:enumerate} with
equal model priors.

$M_2$ and $M_3$ share their three blocks. $M_2$ pools the outer two into one regime, which
is exact because sufficient statistics are additive across conditionally independent
segments, while $M_3$ gives them separate means. Computing the three blocks once and
combining them reduces the work from five passes over the configuration set to three.

Marginal posteriors over each changepoint follow by summation. The excursion indicator is
obtained from a difference array, so the cost is $O(|\mathcal{T}|)$ rather than
$O(|\mathcal{T}| \cdot n)$:
\[
p_{\text{onset}}(t) = \!\!\sum_{(t,\tau_2) \in \mathcal{T}}\!\! W(t,\tau_2),
\qquad
\Pr(h \text{ elevated}) = \!\!\sum_{\tau_1 \leq h < \tau_2}\!\! W(\tau_1,\tau_2),
\qquad
W \propto e^{\ell}.
\]

Reported intervals are the shortest contiguous window carrying $95\%$ of the marginal rather
than a highest-density set. The posterior over a changepoint is often multimodal, and a set
of disjoint hours does not correspond to a date range. The interval is conservative where
modes separate.

\subsection{Result and calibration}

\begin{table}[htbp]
\centering
\small
\begin{tabular}{@{}lrrr@{}}
\toprule
Model & Changepoints & $\log Z$ & Posterior \\
\midrule
$M_0$ no change            & 0 & $605.8$ & $2.7\times10^{-43}$ \\
$M_1$ permanent shift      & 1 & $645.4$ & $4.5\times10^{-26}$ \\
$M_2$ excursion            & 2 & $702.8$ & $0.384$ \\
$M_3$ two free changes     & 2 & $703.3$ & $0.616$ \\
\bottomrule
\end{tabular}
\caption{Marginal likelihood of each segmentation of hourly log dispersion. Posterior mass
on the two-changepoint models is $1.000$ to three decimals. Source:
\code{02\_model\_evidence.csv}.}
\end{table}

\begin{figure}[H]
\centering
\includegraphics[width=\linewidth]{02_changepoint_posterior.png}
\caption{Hourly dispersion with the diurnal profile divided out, and the marginal posterior
over each changepoint. The onset posterior is bimodal, with mass at 24 July and again at the
policy decision on 30 July; the return is better localised. Fitted segment levels are
$0.093$, $0.111$ and $0.091$ pips, a within-hour ratio of $1.19$ to $1.22$ against a pooled
ratio of $3.09$ discussed in Section~\ref{sec:reconcile}. No candidate date was supplied to
any model.}
\end{figure}

A Bayes factor measures how much better $M_2$ fits than $M_0$ on this series. It does not
establish that a series without a changepoint would have looked different. One hundred
replicates were simulated with the sample's length, session structure, fitted $\rho$ and
residual scale, and no change anywhere, then refitted with the same code.

\begin{center}
\small
\begin{tabular}{@{}lr@{}}
\toprule
$\Pr(\text{any change selected} \mid \text{no change})$ & $0.020$ \\
$\Pr(\text{excursion selected} \mid \text{no change})$  & $0.010$ \\
Null $\log$ BF, 95th percentile                         & $-2.8$ \\
Observed $\log$ BF ($M_2$ against $M_0$)                & $+97.1$ \\
Fraction of null replicates the observation exceeds     & $1.000$ \\
\bottomrule
\end{tabular}
\end{center}

The AR(1) term produces that calibration. Under independent errors the same code selects an
excursion on noise in roughly four series out of five. All twenty specifications in
\code{02\_sensitivity.csv} select a two-changepoint model, varying the summary statistic,
bucket width, prior scale, minimum segment length, both admission filters, the
deseasonalisation, the rollover boundary and the error model. None selects $M_0$ or $M_1$.

% =====================================================================
\section{Error correction under a known cointegrating vector}
\label{sec:vecm}

\subsection{Specification}

With the cointegrating vector imposed, the system is a vector error-correction model:
\begin{equation}
\Delta x_t \;=\; c \;+\; \alpha\, \tilde z_{t-1}
\;+\; \sum_{i=1}^{p} \Gamma_i \, \Delta x_{t-i} \;+\; e_t ,
\qquad \alpha \in \mathbb{R}^3,\ \Gamma_i \in \mathbb{R}^{3\times3},
\label{eq:vecm}
\end{equation}
fitted separately inside each estimated regime. A VAR in returns would discard $\tilde z_t$,
which is the only variable in the system that is not a martingale and the only one the
parity condition constrains.

The lag order is $p = 21$ seconds, the cap of the tested grid $\{1,2,3,5,8,13,21\}$. BIC is
still improving at the cap, which is recorded rather than presented as clean selection;
\code{03\_sensitivity.csv} refits at $p \in \{1,5\}$ and the ordering of the regimes is
unchanged.

\subsection{$\lambda = w^{\top}\alpha$ is an identity}

From \eqref{eq:logbasis}, $\tilde z_t = w^{\top}x_t = w^{\top}(x_{t-1} + \Delta x_t)$, so
\begin{equation}
\tilde z_t \;=\; \tilde z_{t-1} + w^{\top} \Delta x_t
\label{eq:zident}
\end{equation}
holds exactly. Pre-multiplying \eqref{eq:vecm} by $w^{\top}$ and substituting,
\[
\Delta \tilde z_t \;=\; w^{\top}c \;+\; \underbrace{(w^{\top}\alpha)}_{\textstyle \lambda}
\tilde z_{t-1} \;+\; \sum_{i=1}^{p} (w^{\top}\Gamma_i)\,\Delta x_{t-i} \;+\; w^{\top}e_t .
\]
$\lambda$ is therefore the one-step closure rate of the basis, obtained from the system
without a separate regression.

\subsection{The sign restriction is a test}

Signs are predicted before estimation. If $\tilde z > 0$ the direct rate is rich against the
synthetic, so closure requires the direct leg to fall and the synthetic to rise:
\[
\alpha_{\AJ} < 0, \qquad \alpha_{\AU} > 0, \qquad \alpha_{\UJ} > 0,
\qquad \text{hence} \qquad \lambda = -\alpha_{\AU} - \alpha_{\UJ} + \alpha_{\AJ} < 0 .
\]
All four hold in all three regimes at $|t|$ between $2.4$ and $11.5$
(\code{04\_weak\_exogeneity.csv}).

\subsection{Two half-lives}

$\lambda$ describes the first second only. The transitory terms can accelerate or delay the
remainder, so the decay of a full dislocation requires the system. Define the state
\[
s_t \;=\; \bigl(\tilde z_t,\; \Delta x_t^{\top},\; \Delta x_{t-1}^{\top},\;\dots,\;
\Delta x_{t-p+1}^{\top}\bigr)^{\top} \in \mathbb{R}^{1+3p}.
\]
Combining \eqref{eq:vecm} with the identity \eqref{eq:zident} gives $s_t = M s_{t-1}$ in the
absence of shocks, where
\begin{equation}
M \;=\;
\begin{pmatrix}
1 + w^{\top}\alpha & w^{\top}\Gamma_1 & \cdots & w^{\top}\Gamma_p \\[2pt]
\alpha & \Gamma_1 & \cdots & \Gamma_p \\[2pt]
0 & I_3 & & \\
& & \ddots & \\
0 & & I_3 & 0
\end{pmatrix}.
\label{eq:companion}
\end{equation}
The first row of \eqref{eq:companion} is the identity \eqref{eq:zident}, not an additional
assumption. The decay path of a unit dislocation is $h(k) = (M^k e_1)_1$, and the
full-system half-life is the first crossing of $\tfrac12$, interpolated between the
bracketing steps.

The spectral radius of $M$ is the stationarity check the error-correction reading requires.
With the cointegrating vector imposed the system carries no unit root in this state, so an
eigenvalue at or above one would mean the regime's basis is not mean-reverting. Observed
values are $0.9889$, $0.9957$ and $0.9841$.

\subsection{Precision is not uniform across $\alpha$}

The three components of $\alpha$ are estimated far less precisely than their combination.
Each leg's own volatility is orders of magnitude larger than the basis, so regressing that
leg's return on the basis carries a large error variance, while $\lambda$ is a combination
in which the common moves cancel by construction. The closure-rate claim is consequently
better supported than the leg-split claim, and the tables do not present them as equally
reliable.

\subsection{HAC for a single linear functional}

Increments carry first-order autocorrelation between $-0.30$ and $-0.42$ from bid-ask
bounce, so OLS standard errors are wrong by a non-negligible factor. Standard errors are
Newey-West.

The full HAC matrix for a three-equation system with $K$ regressors is $3K \times 3K$ and
requires lagged cross-products of an $(n, 3K)$ array. At a million seconds and sixty-five
regressors it is not worth forming, and it is not needed. Only the $\tilde z_{t-1}$
coefficient is of interest, and for a single linear functional the sandwich collapses.
Writing $u_t \in \mathbb{R}^3$ for the residual vector and $e_j$ for the selector of the
error-correction column,
\[
\hat\alpha - \alpha
\;=\; \sum_t \underbrace{e_j^{\top}(X^{\top}X)^{-1} x_t}_{\textstyle h_t \in \mathbb{R}}
\; u_t^{\top},
\]
so the required $3 \times 3$ covariance is the long-run covariance of the single vector
series $\varphi_t = h_t u_t$.

That covariance is computed from overlapping block sums rather than a loop over lags. Let
$B_s = \sum_{i=0}^{m-1} \varphi_{s+i}$. Expanding,
\[
\frac{1}{m}\sum_s B_s B_s^{\top}
= \frac{1}{m}\sum_s \sum_{i,j=0}^{m-1} \varphi_{s+i}\varphi_{s+j}^{\top}
= \sum_{|k| < m} \frac{m - |k|}{m}\, \Gamma_k
= \sum_{|k| < m}\left(1 - \frac{|k|}{m}\right)\Gamma_k ,
\]
since exactly $m - |k|$ ordered pairs $(i,j)$ in $[0,m)^2$ satisfy $i - j = k$. This
reproduces the Bartlett kernel with truncation $m-1$ exactly, up to the first and last $m-1$
rows of each contiguous run. With runs hours long and $m$ of the order of a minute, that
edge term is around $10^{-5}$ of the total. Blocks that would straddle a break are refused.

\subsection{Row admission}

A lag window is admitted only where the whole span of $p+2$ consecutive seconds is
contiguous in clock time, outside the daily rollover window, inside one regime and one
trading day, and labelled with regime posterior at least $90\%$. Nothing is differenced
across a weekend, a closure, a rollover exclusion or a regime boundary.

Using the posterior rather than the hard boundary places a buffer around each changepoint
that is derived from the data instead of chosen. Both changepoints span market closures, so
some buffer is required. The row set is fixed once at the lag cap and reused for every lag
order and every regime, so no two numbers in the output are computed on different samples.

The segment boundaries come from the median segmentation of Section~\ref{sec:cp} and are
read from the labels rather than typed as dates. The onset boundary is a choice between two
posterior modes and is discussed in Section~\ref{sec:onset}; the alternative is carried as a
row of \code{03\_sensitivity.csv} so its cost can be read off.

One further choice matters for Section~\ref{sec:reconcile}. The error-correction term is
centred on the day. Centring on the regime leaves drift in the daily mean inside the residual, where
the estimator reads it as persistence: that variant returns a stress half-life of $967$
seconds against a baseline of $155$.

% =====================================================================
\section{Attribution}
\label{sec:attrib}

Three routes, sharing no machinery, in increasing order of what each assumes.

\subsection{By identity}

From \eqref{eq:zident}, the change in the basis over any window is exactly the sum of three
leg contributions:
\[
\tilde z_{t_1} - \tilde z_{t_0} \;=\; \sum_{k \in \{\AU,\UJ,\AJ\}} w_k
\bigl(x^{k}_{t_1} - x^{k}_{t_0}\bigr).
\]
Applying this to the opening and closing window of each detected episode splits both by
arithmetic, with no model and no lag order.

\begin{table}[htbp]
\centering
\small
\begin{tabular}{@{}lrrr@{}}
\toprule
Share of the move & pre & stress & post \\
\midrule
Opening: AUD/USD $+$ USD/JPY & $0.86$ & $1.12$ & $1.05$ \\
Closing: AUD/JPY             & $0.64$ & $0.71$ & $0.48$ \\
\midrule
Median seconds to open  & 2 & 19 & 4 \\
Median seconds to close & 3 & 24 & 4 \\
\bottomrule
\end{tabular}
\caption{Episode-level attribution. Shares can exceed one because leg contributions can
oppose. Source: \code{04\_attribution.csv}.}
\end{table}

\subsection{By correlation and by test}

Cross-correlations between each leg's increments and the basis's are computed at matched
leads and lags out to a minute, on pairs inside one unbroken run. Block Granger tests are
then run inside \eqref{eq:vecm}.

A degeneracy has to be handled. The basis is a fixed combination of the three legs, so at
every lag the three legs' increments together span the basis's own lagged increment.
Dropping any one leg's block removes part of the basis's own autoregression along with that
leg's information, and when the basis is persistent all three blocks then test large for
that reason alone. The headline test therefore conditions on the basis's own history
explicitly: for each leg, a smaller model containing the basis level, the basis's own lagged
increments and that one leg's lagged increments, asking what the leg adds beyond the basis's
own memory. The full-system figure is reported beside it as a control. Neither model is
refitted; both are column subsets of one extended design whose cross-products are
accumulated once.

\subsection{Effect sizes rather than $p$-values}

Every test runs on between one hundred thousand and a million seconds. At that size a
coefficient of no economic consequence rejects at any conventional level. Each test
therefore carries an effect size, the incremental predictable move in pips per second,
beside a placebo threshold computed from the same statistic with that leg's increments drawn
from a different day. The placebo preserves the leg's own autocorrelation and volatility
clustering and destroys only the cross-leg timing, so it is the distribution of the
statistic when the null holds and everything else is unchanged.

The results separate cleanly. Cross-leg effects clear their placebo by $3.5$ to $10.6$
times. Leg-to-basis effects, conditioned on the basis's own memory, explain between $0.01\%$
and $0.07\%$ of basis increment variance; AUD/JPY fails Holm correction in all three regimes
and USD/JPY in two. The basis is close to unforecastable from its own components while those
components remain highly forecastable from each other.

\subsection{Why no Hasbrouck or Gonzalo-Granger shares}

Those measures are defined for two prices of one asset: one cointegrating vector, one common
trend, and loadings that split into shares summing to one. This system has three prices and
one restriction, so it has two common trends, an AUD trend and a JPY trend, and
$\alpha_{\perp}$ is $3 \times 2$ rather than a vector. No scalar share per leg is identified
without a normalisation the data does not supply. The weak-exogeneity tests are the part of
that framework this system supports, and they are what is reported.

% =====================================================================
\section{The transaction-cost band}
\label{sec:band}

\subsection{The two round trips}

Sections~\ref{sec:cp} and \ref{sec:vecm} work on mid quotes, which is the correct series for
measuring how prices adjust and the wrong one for asking whether a dislocation was
executable. A round trip crosses three spreads. Writing $b$ and $a$ for bid and ask in
AUD/JPY price units,
\[
e_A = \AJ_b - \AU_a \UJ_a \quad (\text{sell direct, buy synthetic}), \qquad
e_B = \AU_b \UJ_b - \AJ_a \quad (\text{buy direct, sell synthetic}).
\]
Each is the profit per unit of AUD from one complete circuit at prices that were quoted.
They cannot be positive together, since $e_A > 0$ and $e_B > 0$ would require
$\AJ_b > \AJ_a$.

\subsection{The band, and the exact cancellation}

Let $z$ be the mid basis. Then $e_A = z - c_A$ and $e_B = -z - c_B$ with
\[
c_A = (\AJ_m - \AJ_b) + (\AU_a \UJ_a - \AU_m \UJ_m), \qquad
c_B = (\AJ_a - \AJ_m) + (\AU_m \UJ_m - \AU_b \UJ_b),
\]
so nothing is executable while $z \in [-c_B, +c_A]$. The width of that band is
\[
W = c_A + c_B = \underbrace{(\AJ_a - \AJ_b)}_{s_{\AJ}}
\;+\; \bigl(\AU_a \UJ_a - \AU_b \UJ_b\bigr).
\]
Writing each quote as its mid plus or minus half a spread, the second term expands as
\[
\left(\AU_m + \tfrac{s_{\AU}}{2}\right)\!\left(\UJ_m + \tfrac{s_{\UJ}}{2}\right)
- \left(\AU_m - \tfrac{s_{\AU}}{2}\right)\!\left(\UJ_m - \tfrac{s_{\UJ}}{2}\right)
= \AU_m s_{\UJ} + \UJ_m s_{\AU},
\]
the quadratic terms cancelling in the difference. Therefore
\begin{equation}
\boxed{\;W \;=\; s_{\AJ} \;+\; \AU_m \, s_{\UJ} \;+\; \UJ_m \, s_{\AU}\;}
\label{eq:band}
\end{equation}
holds exactly rather than to first order. The band decomposes additively across the three
legs with no interaction term, so which leg sets the bound is determined by arithmetic.
Nothing in \eqref{eq:band} is estimated and no cost parameter is assumed.

\subsection{The band in units of the basis}

Section~\ref{sec:cp} measured a numerator. Equation \eqref{eq:band} supplies the
denominator. The quantity of interest is $W / \operatorname{MAD}(z)$ per regime, computed on
the same seconds.

\begin{table}[htbp]
\centering
\small
\begin{tabular}{@{}lrrr@{}}
\toprule
 & pre & stress & post \\
\midrule
Band width $W$, median (pips)      & $2.61$ & $3.35$ & $2.57$ \\
Basis dispersion, MAD (pips)       & $0.112$ & $0.264$ & $0.095$ \\
Band / dispersion                  & $23.3$ & $12.7$ & $27.1$ \\
\bottomrule
\end{tabular}
\caption{Rows restricted to seconds where all three legs quoted within one second and the
regime posterior is at least $90\%$. Source: \code{05\_execution.csv}.}
\end{table}

Both quantities widened. The dispersion grew by $2.35$ and the band by $1.29$, so the
band-to-gap ratio approximately halved. The basis moved closer to the band and remained an
order of magnitude inside it.

\begin{figure}[H]
\centering
\includegraphics[width=\linewidth]{05_bands.png}
\caption{The mid basis against the band it would have had to cross. On the log axis of the
lower panel a parallel shift indicates that the band-to-gap ratio did not change.}
\end{figure}

\subsection{Sensitivity to the spread level}

Both vendors publish aggregator quotes, which are wider than interdealer, so the comparison
is biased toward finding no executable dislocation. Every spread is therefore scaled by a
factor $k$, equivalent to $e_A(k) = z - k c_A$ and $e_B(k) = -z - k c_B$ with the mid basis
held fixed, and the share of seconds with a positive edge is reported as a function of $k$.
Halving every quoted spread leaves $99.95\%$ of calm seconds inside the band; only at a
tenth of quoted do the calm regimes reach $22.6\%$ and $14.4\%$ (\code{05\_haircut.csv}).

\subsection{Why a positive edge is not an opportunity}

The grid is one second wide and the three legs are separate feeds, so a row records the last
print of each leg inside that second rather than three prices known to have stood together.
Triangular parity requires simultaneity across all three. A positive edge computed from this
grid is therefore an upper bound on what was available and not a measurement of it: when
prices travel far enough inside a second, the separation between prints produces an apparent
edge, and second resolution cannot distinguish the two readings. No finer data exists for
these pairs from either vendor.

Two diagnostics are reported in place of an assertion. A real opportunity decays as the
delay between signal and fill grows and a measurement artefact does not; the observed count
is flat across execution delays of zero to five seconds. A displacement that was present is
indifferent to how fast the market was moving and non-simultaneity is not; the count is not
concentrated in the fastest quintiles of within-second travel. Both point away from reading
the stress count as opportunity.

% =====================================================================
\section{Reconciliation}
\label{sec:reconcile}

Five routes to the closure speed of the basis are distributed across the preceding sections
and they do not agree.

\begin{table}[htbp]
\centering
\small
\begin{tabular}{@{}llr@{}}
\toprule
Route & What it estimates & Stress / calmest \\
\midrule
Episode ruler            & time for a large excursion to halve & $9.3$ \\
AR(1) on the level       & decay of an average deviation, one parameter & $9.7$ \\
VECM, one-step           & decay implied by $\lambda$ alone & $2.9$ \\
VECM, full system        & decay of the 21-lag system run down & $4.1$ \\
VECM, day by day, median & the typical day rather than the typical second & $1.4$ \\
\bottomrule
\end{tabular}
\caption{Source: \code{06\_closure\_reconciliation.csv}.}
\end{table}

\begin{figure}[H]
\centering
\includegraphics[width=\linewidth]{03_closure_speed.png}
\caption{Each day fitted on its own, with no regime label supplied. The pooled stress mean
of $155$ seconds sits above every stress day but one, which is the disagreement the rest of
this section accounts for. Discussed in Section~7.4.}
\end{figure}

\subsection{Decomposing the level change}

For a stationary AR(1) level with innovation scale $\sigma$ and coefficient $\rho$,
\begin{equation}
\operatorname{sd}(z) \;=\; \frac{\sigma}{\sqrt{1-\rho^{2}}},
\qquad\text{so}\qquad
\log \frac{\operatorname{sd}(z)_A}{\operatorname{sd}(z)_B}
= \underbrace{\log \frac{\sigma_A}{\sigma_B}}_{\text{noise}}
+ \underbrace{\log \frac{\sqrt{1-\rho_B^{2}}}{\sqrt{1-\rho_A^{2}}}}_{\text{persistence}} .
\label{eq:ampl}
\end{equation}
The split is additive on the log scale. The innovation scale of the basis does not need
separate estimation, because $\varepsilon_z = w^{\top}e$ and therefore
\begin{equation}
\sigma^2 \;=\; w^{\top} \Sigma\, w
\;=\; \sigma_{\AU}^2 + \sigma_{\UJ}^2 + \sigma_{\AJ}^2
+ 2\sigma_{\AU}\sigma_{\UJ}\rho_{\AU\UJ}
- 2\sigma_{\AU}\sigma_{\AJ}\rho_{\AU\AJ}
- 2\sigma_{\UJ}\sigma_{\AJ}\rho_{\UJ\AJ},
\label{eq:sigmaz}
\end{equation}
reconstructible from the per-leg residual scales and correlations that \eqref{eq:vecm}
already produces. Substituting them reproduces the measured level dispersion to within
$6\%$ to $17\%$ across all three regimes, which confirms that \eqref{eq:ampl} and
\eqref{eq:sigmaz} describe the same object.

\subsection{What the pooled statistic measures}

Equation \eqref{eq:ampl} attributes the widening to larger shocks and slower reversion. That
attribution is incomplete, because $\operatorname{sd}(z)$ is a pooled quantity. Splitting the
level variance on the trading day $d$,
\begin{equation}
\operatorname{Var}(z)
\;=\; \underbrace{\mathbb{E}_d\bigl[\operatorname{Var}(z \mid d)\bigr]}_{\text{within-day}}
\;+\; \underbrace{\operatorname{Var}_d\bigl(\mathbb{E}[z \mid d]\bigr)}_{\text{day-to-day}} .
\label{eq:totalvar}
\end{equation}
Both terms are evaluated on the same seconds by \code{07\_boundary\_sweep.py}, which
accumulates per-hour sums once and aggregates them into per-day cells. The same widening
then reads differently at four nested scopes, and nothing about the market changes between
the rows.

\begin{table}[htbp]
\centering
\small
\begin{tabular}{@{}lr@{}}
\toprule
Scope of the measurement & Stress / calm \\
\midrule
Within the hour (MAD)                          & $1.21$ \\
Within the day (sd)                            & $1.98$ \\
Pooled across the window, robust (MAD)         & $2.91$ \\
Pooled across the window, non-robust (sd)      & $4.22$ \\
\midrule
Day-to-day share of level variance, stress     & $80.6\%$ \\
Day-to-day share of level variance, calm       & $11.6\%$ \\
\bottomrule
\end{tabular}
\caption{Applying \eqref{eq:totalvar} at the segmentation in use. Calm pools the pre and
post segments. Source: \code{07\_boundary\_named.csv}.}
\end{table}

Four fifths of what the pooled statistic calls dispersion under stress is the level moving
between days. Each step up the ladder admits another layer of that movement. The within-hour
figure of $1.21$ agrees with the $1.19$ to $1.22$ the changepoint model returned
independently from hourly buckets, which is the check that the decomposition and the fit
describe the same object.

\subsection{The decomposition does not turn on the regime boundary}
\label{sec:sweep}

Section~\ref{sec:onset} records that the onset posterior is bimodal and that the boundary in
use is the median rather than the mode. Every quantity above is conditioned on that choice.

Rather than evaluate the decomposition at one boundary, or at two,
\code{07\_boundary\_sweep.py} evaluates it at every hourly onset between 23 and 31 July under
both candidate returns, $128$ boundaries in total. The decomposition is a closed form in
per-day sufficient statistics, so a boundary costs $O(\text{hours})$ rather than
$O(\text{seconds})$. Nothing is refitted; the basis series is re-aggregated.

\begin{table}[htbp]
\centering
\small
\begin{tabular}{@{}lrrr@{}}
\toprule
 & median rule (in use) & modal rule & all $128$ boundaries \\
\midrule
Onset                        & 30 Jul 15:00 & 24 Jul 17:00 & 23 to 31 Jul \\
Return                       & 09 Aug 10:00 & 13 Aug 19:00 & both \\
Stress window (hours)        & $183$ & $331$ & $160$ to $354$ \\
Day-to-day share, stress     & $80.6\%$ & $78.7\%$ & $78.7\%$ to $87.3\%$ \\
Day-to-day share, calm       & $11.6\%$ & $11.8\%$ & $10.2\%$ to $12.9\%$ \\
Within-hour MAD ratio        & $1.21$ & $1.13$ & $1.12$ to $1.21$ \\
Pooled MAD ratio             & $2.91$ & $1.63$ & \\
\bottomrule
\end{tabular}
\caption{The decomposition as a function of the boundary. Source:
\code{07\_boundary\_named.csv} and \code{07\_boundary\_sweep.csv}.}
\end{table}

\begin{figure}[H]
\centering
\includegraphics[width=\linewidth]{07_boundary_sweep.png}
\caption{The decomposition at every candidate onset, under both candidate returns. The
upper panel is flat across the swept range, so the headline does not turn on the boundary.
In the lower panel the pooled ratio rises with the boundary while the two decomposed
ratios do not, and the gap between them is the pooling effect.}
\end{figure}

The day-to-day share stays between $78.7\%$ and $87.3\%$ across every boundary tested,
against $10.2\%$ to $12.9\%$ in the calm segments, over stress windows ranging from $160$ to
$354$ hours. The within-hour widening moves by $8\%$ across the whole sweep, which makes it
the most stable quantity in the project.

The pooled MAD ratio is the exception, falling from $2.91$ to $1.63$ when the boundary moves
to the modal onset. The quantity this report treats as the headline is stable under the
boundary choice; the quantity it argues against reporting is the one that is not.

\subsection{Effect on the persistence estimates}

Write $z_t = \mu_d + u_t$ with $\mu_d$ constant within a day and $u_t$ an AR(1) with
coefficient $\rho_w$. Almost every lag-one pair lies inside a day, so
\[
\rho_{\text{pooled}}
= \frac{\sigma_{\text{day}}^{2} + \rho_w \sigma_{\text{within}}^{2}}
       {\sigma_{\text{day}}^{2} + \sigma_{\text{within}}^{2}} .
\]
Inverting at the shares above:

\begin{center}
\small
\begin{tabular}{@{}lrrr@{}}
\toprule
 & pre & stress & post \\
\midrule
Pooled AR(1), as reported        & $0.851$ & $0.976$ & $0.792$ \\
Implied within-day AR(1)         & $0.825$ & $0.723$ & $0.786$ \\
Implied within-day half-life (s) & $3.6$ & $2.1$ & $2.9$ \\
\bottomrule
\end{tabular}
\end{center}

The stress regime has the lowest within-day persistence of the three. Its pooled AR(1) of
$0.976$ is almost entirely the daily-mean component.

This is a two-component calculation on published values rather than a re-estimation. It is
corroborated by a route sharing none of its machinery: the VECM refitted day by day removes
the daily level by construction and returns median half-lives of $43.6$, $38.4$ and $27.6$
seconds, placing the stress regime between the other two.

Three estimators that remove the daily level return a small effect. Three that retain it
return a large one. The choice of estimand moves the reported figure by a factor of seven.

\subsection{One day, not one regime}

The daily fits in Figure~3 resolve what remains. Eight of the nine stress days fit at
$14$ and $67$ seconds, inside the pre-stress range of $10$ to $144$. The ninth is 31 July,
the session following the announcement at 30 July 23:30 in the data's clock, at $1{,}003$
seconds: seven times slower than the next slowest day
in the sample and fifteen times slower than the next slowest stress day. The pooled stress
mean of $155$ seconds is above every stress day except that one.


Across the conventions in \code{03\_sensitivity.csv} the ratio spans $1.5$ to $24.8$. The
direction is stable across all of them; the magnitude is not determined by the data alone.
The supportable statement is that one day was extraordinary.

% =====================================================================
\section{Limitations}
\label{sec:limits}

\subsection{Two days fail the data-quality gate}

The daily-bias check of Section~\ref{sec:grid} rejects seven of $53$ days at a threshold of
$0.10$ pips. Two of them by a factor of fifteen.

\begin{table}[htbp]
\centering
\small
\begin{tabular}{@{}lrrrr@{}}
\toprule
Day & AUD/USD & USD/JPY & AUD/JPY & Open-market seconds \\
\midrule
11 Jul & $+0.143$ & $+0.338$ & $-0.228$ & $86{,}400$ \\
12 Jul & $+0.137$ & $+0.320$ & $-0.217$ & $61{,}200$ \\
31 Jul & $-0.045$ & $-0.104$ & $+0.068$ & $86{,}400$ \\
1 Aug  & $-0.977$ & $-2.247$ & $+1.464$ & $86{,}400$ \\
2 Aug  & $-0.991$ & $-2.252$ & $+1.468$ & $61{,}200$ \\
8 Aug  & $+0.176$ & $+0.393$ & $-0.258$ & $86{,}400$ \\
9 Aug  & $+0.180$ & $+0.402$ & $-0.265$ & $61{,}200$ \\
\bottomrule
\end{tabular}
\caption{Median signed gap between each leg and the value the other two imply, in that leg's
own pips. Source: \code{00\_daily\_bias.csv}.}
\end{table}

The gate returns a non-zero exit code on this sample. It exists to catch a persistent
sub-pip bias, which the tail statistic cannot see: the $99.9$th percentile of the absolute
gap passes at $3.21$, $7.29$ and $4.83$ pips against a $25$-pip threshold.

The triangle cannot attribute the discrepancy to a leg. The three columns are one
discrepancy expressed three ways and any single leg reproduces it exactly. Attribution
requires a reference outside the triangle, and none is available: Dukascopy carries the same
level on those days, but HistData sources from Dukascopy, so agreement between them
establishes that the level is present in a shared upstream feed rather than that it is real.
Settling it requires a feed from a different liquidity pool. The offset is documented and
carried, not fixed.

1 and 2 August carry a standing basis offset near $-1.47$ pips across $147{,}600$
open-market seconds, which is $22\%$ to $24\%$ of the stress regime depending on the
admission rule. The implied contribution to the regime mean basis is $-0.33$ to $-0.35$
pips. The measured stress mean basis is $-0.27$ pips, against $+0.016$ pre and $+0.005$
post. The two days account for the entire mean level shift attributed to the stress regime.

Estimates insensitive to this: the changepoint dates, because a level offset constant within
a day leaves within-hour dispersion unchanged and neither estimated changepoint falls near
1 or 2 August; and the error-correction coefficients, because the VECM centres on the day.

Estimates sensitive to it: the regime-level standard deviation, the pooled AR(1) of the
level, and the positive-edge counts of Section~\ref{sec:band}.

\subsection{The onset is bimodal, and the regime boundary is a choice}
\label{sec:onset}

The onset posterior has two modes. One sits at 24 July 17:00, the other at the policy
decision on 30 July, and the $95\%$ window between them spans a weekend closure. Across the
$17$ same-feature specifications of \code{02\_sensitivity.csv} the split is $11$ at 24 July,
$4$ at 30 July, $1$ at 29 July and $1$ at 1 July.

The boundary this project uses is 30 July, which is not the date the estimator selects most
often. There are two reasons and they are of different kinds.

The first is methodological. The per-second labels written by
\code{02\_changepoint\_detection.py} are the \emph{median} segmentation: a bucket is called
elevated when more than half the posterior mass says it is. 30 July is what that rule
returns. 24 July is the posterior mode, a taller single spike carrying less than half the
total mass. Reporting the mode while labelling on the median would leave the headline
describing a different window from the one every downstream script consumes.

The second is historical, and was a researcher degree of freedom rather than a result. The
project was originally scoped around the 31 July policy decision, so the boundary sitting
$8.5$ hours before that announcement is the one it was built to find. That is a prior, and
Section~\ref{sec:sweep} replaces it with a measurement: the decomposition is evaluated at
every candidate boundary rather than at the chosen one.

The cost of the alternative is measured rather than argued. Taking the 24 July onset moves
$338{,}396$ seconds from pre into stress and the closure ratio falls from $4.06$ to $3.35$
(\code{03\_sensitivity.csv}, row \emph{02's modal onset, 24 Jul 17:00}). The chosen boundary
is therefore the one under which the estimated shift is larger, and a reader who prefers
the modal onset should read $3.35$ and adjust Section~\ref{sec:reconcile} accordingly.

One observation bears on which mode is the more credible without settling it. The
24 July 17:00 bucket covers 17:30 to 17:59:59, and it contains two of the twenty widest
episodes in the sample: $7.335$ pips at 17:44:08, the widest in the entire pre-stress
segment, and $5.65$ pips at 17:58:14. Both fall in the half hour after the rollover filter
reopens, which the boundary audit of the next subsection shows carries $21.5$ times the
expected episode rate. Part of the 24 July mode may therefore be the filter boundary rather
than the market. That is not resolved here.

A change falling inside a closure can only be recorded at the reopen, so the tables report
the bracket the data supports rather than a modal hour. The return is better behaved,
clustering at 12 to 13 August in $14$ of the $17$ same-feature specifications.

\subsection{Censoring at the filter boundary}

Episode detection excludes 17:00 to 17:30, so the detector sees a session opening at
17:30:00 and closing at 16:59:59. Nine of the twenty widest episodes start in hour 17, which
holds $2.1\%$ of admissible seconds, at $21.5$ times the expected rate. Two begin exactly at
17:30:00 and three end exactly at 16:59:59, against an expectation of $0.0002$ for each.
These are censored observations rather than starts and ends, and the widest episode in the
sample, $8.911$ pips on 4 August, is one of them. Moving the boundary to 18:00 leaves the
selected segmentation unchanged.

\subsection{Scope}

One event, two months, one currency triangle. The results describe a single episode and are
not distributional claims about FX stress in general. The quotes are retail aggregator
quotes, which biases Section~\ref{sec:band} toward finding no executable dislocation; the
haircut ladder quantifies that bias rather than removing it.

\subsection{Summary}

Supported by the evidence:

\begin{itemize}[leftmargin=1.6em,itemsep=3pt]
  \item Hourly dispersion is best described by a bounded excursion that returns rather than
  a permanent shift.
  \item Within-hour dispersion rose by $1.21$ times, and that figure is stable across all
  $128$ candidate regime boundaries.
  \item The transaction-cost band exceeded the median basis by $12.7$ times at its tightest.
  \item The two dollar legs opened the observed dislocations and the cross leg closed them.
  \item The basis is close to unforecastable from its own components.
  \item 31 July had a closure half-life of $1{,}003$ seconds.
\end{itemize}

Not supported: a permanent regime change, a uniformly slower stress regime, a threefold
widening of the intraday basis, or any executable dislocation.

% =====================================================================
\section*{Appendix: where each claim is written down}

\small
\begin{tabularx}{\linewidth}{@{}lX@{}}
\toprule
\textbf{File} & \textbf{Claim it supports} \\
\midrule
\code{00\_daily\_bias.csv} & Seven days breach the bias gate; 1 and 2 August by fifteenfold \\
\code{00\_reconciliation.csv} & Tail agreement of each leg with its implied value \\
\code{01\_summary\_stats.csv} & Moments of the basis; excess kurtosis $277$ \\
\code{01\_closures.csv} & Eight closures detected from the data \\
\code{02\_model\_evidence.csv} & Posterior over the four segmentations \\
\code{02\_null\_calibration.csv} & Behaviour of the estimator when no change exists \\
\code{02\_sensitivity.csv} & Twenty specifications, all selecting two changepoints \\
\code{03\_var\_fit.csv} & Residual scales, correlations, sd of the daily mean basis \\
\code{03\_error\_correction.csv} & $\alpha$, $\lambda$, HAC intervals, both half-lives \\
\code{03\_closure\_speed.csv} & Daily $\lambda$ and half-life, 53 sessions \\
\code{03\_sensitivity.csv} & Closure ratio across conventions, $1.5$ to $24.8$ \\
\code{04\_attribution.csv} & Opening and closing shares by identity \\
\code{04\_weak\_exogeneity.csv} & Sign restriction on $\alpha$, all three regimes \\
\code{04\_basis\_causality.csv} & Effect sizes against placebo thresholds \\
\code{05\_execution.csv} & Band width, dispersion, band-to-gap ratio \\
\code{05\_haircut.csv} & Share clearing as a function of the spread scale \\
\code{05\_latency.csv}, \code{05\_edge\_velocity.csv} & Diagnostics on the simultaneity limit \\
\code{06\_normalisation.csv} & Noise and persistence split of the level change \\
\code{06\_closure\_reconciliation.csv} & Five estimands of closure speed \\
\code{06\_boundary\_audit.csv} & Censoring at the rollover boundary \\
\code{07\_boundary\_named.csv} & The decomposition under each summarisation of the onset \\
\code{07\_boundary\_sweep.csv} & The decomposition at all $128$ candidate boundaries \\
\midrule
\code{02\_changepoint\_posterior.png} & Excursion, bimodal onset, fitted segment levels \\
\code{03\_closure\_speed.png} & Pooled stress mean set by one day \\
\code{05\_bands.png} & Band against basis on a log axis \\
\code{07\_boundary\_sweep.png} & The decomposition against the regime boundary \\
\bottomrule
\end{tabularx}

\end{document}